% GENERATED FILE. Edit canonical lessons, then run npm run generate:books.
% canonical-source-hash: fnv1a64:f37a8c7830499387
% canonical-lessons: ML-C365-parati, ML-C365-kurravali, ML-C365-anvesanam, ML-C365-arasrru-ceyyuka, ML-C365-piticcetukkuka

\chapter{At the Police Station}
\label{ch:ml-at-the-police-station}
\hlchaptermodality{365}

\begin{chapteropening}
\textbf{By the end of this chapter:} \emph{I can say a complaint, a criminal, an investigation, to arrest and to seize.}

Complete the last lesson of chapter 365: I can say a complaint, a criminal, an investigation, to arrest and to seize.
\end{chapteropening}

\section[\texorpdfstring{\ml{പരാതി} (parāti)}{പരാതി (parāti)}]{\ml{പരാതി} (parāti) — a complaint}

\label{lesson:ML-C365-parati}



\begin{quote}
Before the new one: say the Malayalam for to steal, then the Malayalam for to break open.
\end{quote}

\subsection*{You'll want to know: \ml{പരാതി}}
\textbf{\ml{പരാതി}} — \emph{parāti} — \textquotedblleft{}a complaint\textquotedblright{}.

\begin{quote}
\ml{അച്ഛൻ} \ml{പോലീസിൽ} \ml{ഒരു} \ml{പരാതി} \ml{കൊടുത്തു}.
\end{quote}

\begin{quote}
\emph{acchan pōlīsil oru parāti koṭuttu.}
\end{quote}

\begin{quote}
Father filed a complaint with the police.
\end{quote}

A complaint is given: \ml{പരാതി} \ml{കൊടുക്കുക}. \ml{പരാതിപ്പെടുക} is the verb.

\begin{cousinweb}
The origin of \textbf{\ml{പരാതി}} (\emph{parāti}) is not settled. It is the noun behind \ml{പരാതിപ്പെടുക}, \textquotedblleft{}to complain\textquotedblright{}, and the usual word for a complaint to the police or an office.
\end{cousinweb}

\begin{grammarlens}[title={What you've built}]
One more word for telling what happened.
\end{grammarlens}

\subsection*{Your turn}
\begin{itemize}
\raggedright
  \item \emph{Say it:} \emph{parāti}
  \item \emph{Say it:} \emph{parāti}, once more
  \item \emph{Say it:} \emph{acchan pōlīsil oru parāti koṭuttu.} — Father filed a complaint with the police.
  \item \emph{Recall:} say the Malayalam for to steal, then the Malayalam for to break open, then say \emph{parāti} again
\end{itemize}

\begin{tcolorbox}[breakable,colback=teal!4,colframe=teal!35!black,title={Before you move on}]
What does \ml{പരാതി} mean? (\textquotedblleft{}A complaint\textquotedblright{}.) Say it once more.
\end{tcolorbox}
\section[\texorpdfstring{\ml{കുറ്റവാളി} (kuṟṟavāḷi)}{കുറ്റവാളി (kuṟṟavāḷi)}]{\ml{കുറ്റവാളി} (kuṟṟavāḷi) — a criminal; the culprit}

\label{lesson:ML-C365-kurravali}



\begin{quote}
Before the new one: say the Malayalam for to break open, then the Malayalam for a complaint.
\end{quote}

\subsection*{You'll want to know: \ml{കുറ്റവാളി}}
\textbf{\ml{കുറ്റവാളി}} — \emph{kuṟṟavāḷi} — \textquotedblleft{}a criminal; the culprit\textquotedblright{}.

\begin{quote}
\ml{പോലീസ്} \ml{കുറ്റവാളിയെ} \ml{കണ്ടെത്തി}.
\end{quote}

\begin{quote}
\emph{pōlīsŭ kuṟṟavāḷiye kaṇṭetti.}
\end{quote}

\begin{quote}
The police found the culprit.
\end{quote}

\ml{കുറ്റം}, an offence, + -\ml{വാളി}, one who has it.

\begin{cousinweb}
\textbf{\ml{കുറ്റവാളി}} (\emph{kuṟṟavāḷi}) is native Dravidian: \ml{കുറ്റം} (\emph{kuṟṟaṁ}), \textquotedblleft{}a fault, an offence\textquotedblright{}, which DEDR 1851 sets with a root meaning short, lacking, and the ending -\ml{വാളി}, \textquotedblleft{}one who has it\textquotedblright{}.
\end{cousinweb}

\begin{grammarlens}[title={What you've built}]
One more word for telling what happened.
\end{grammarlens}

\subsection*{Your turn}
\begin{itemize}
\raggedright
  \item \emph{Say it:} \emph{kuṟṟavāḷi}
  \item \emph{Say it:} \emph{kuṟṟavāḷi}, once more
  \item \emph{Say it:} \emph{pōlīsŭ kuṟṟavāḷiye kaṇṭetti.} — The police found the culprit.
  \item \emph{Recall:} say the Malayalam for to break open, then the Malayalam for a complaint, then say \emph{kuṟṟavāḷi} again
\end{itemize}

\begin{tcolorbox}[breakable,colback=teal!4,colframe=teal!35!black,title={Before you move on}]
What does \ml{കുറ്റവാളി} mean? (\textquotedblleft{}A criminal; the culprit\textquotedblright{}.) Say it once more.
\end{tcolorbox}
\section[\texorpdfstring{\ml{അന്വേഷണം} (anvēṣaṇaṁ)}{അന്വേഷണം (anvēṣaṇaṁ)}]{\ml{അന്വേഷണം} (anvēṣaṇaṁ) — an investigation; an inquiry}

\label{lesson:ML-C365-anvesanam}



\begin{quote}
Before the new one: say the Malayalam for a complaint, then the Malayalam for a criminal.
\end{quote}

\subsection*{You'll want to know: \ml{അന്വേഷണം}}
\textbf{\ml{അന്വേഷണം}} — \emph{anvēṣaṇaṁ} — \textquotedblleft{}an investigation; an inquiry\textquotedblright{}.

\begin{quote}
\ml{പോലീസ്} \ml{ആ} \ml{മോഷണത്തെക്കുറിച്ച്} \ml{അന്വേഷണം} \ml{തുടങ്ങി}.
\end{quote}

\begin{quote}
\emph{pōlīsŭ ā mōṣaṇattekkuṟiccŭ anvēṣaṇaṁ tuṭaṅṅi.}
\end{quote}

\begin{quote}
The police started an investigation into the theft.
\end{quote}

An investigation is held: \ml{അന്വേഷണം} \ml{നടത്തുക}.

\begin{cousinweb}
\textbf{\ml{അന്വേഷണം}} (\emph{anvēṣaṇaṁ}) is borrowed from Sanskrit \emph{anveṣaṇa}, \textquotedblleft{}seeking for, searching, investigating\textquotedblright{} (Monier-Williams).
\end{cousinweb}

\begin{grammarlens}[title={What you've built}]
One more word for telling what happened.
\end{grammarlens}

\subsection*{Your turn}
\begin{itemize}
\raggedright
  \item \emph{Say it:} \emph{anvēṣaṇaṁ}
  \item \emph{Say it:} \emph{anvēṣaṇaṁ}, once more
  \item \emph{Say it:} \emph{pōlīsŭ ā mōṣaṇattekkuṟiccŭ anvēṣaṇaṁ tuṭaṅṅi.} — The police started an investigation into the theft.
  \item \emph{Recall:} say the Malayalam for a complaint, then the Malayalam for a criminal, then say \emph{anvēṣaṇaṁ} again
\end{itemize}

\begin{tcolorbox}[breakable,colback=teal!4,colframe=teal!35!black,title={Before you move on}]
What does \ml{അന്വേഷണം} mean? (\textquotedblleft{}An investigation; an inquiry\textquotedblright{}.) Say it once more.
\end{tcolorbox}
\section[\texorpdfstring{\ml{അറസ്റ്റ്} \ml{ചെയ്യുക} (aṟasṟṟŭ ceyyuka)}{അറസ്റ്റ് ചെയ്യുക (aṟasṟṟŭ ceyyuka)}]{\ml{അറസ്റ്റ്} \ml{ചെയ്യുക} (aṟasṟṟŭ ceyyuka) — to arrest}

\label{lesson:ML-C365-arasrru-ceyyuka}



\begin{quote}
Before the new one: say the Malayalam for a criminal, then the Malayalam for an investigation.
\end{quote}

\subsection*{You'll want to know: \ml{അറസ്റ്റ്} \ml{ചെയ്യുക}}
\textbf{\ml{അറസ്റ്റ്} \ml{ചെയ്യുക}} — \emph{aṟasṟṟŭ ceyyuka} — \textquotedblleft{}to arrest\textquotedblright{}.

\begin{quote}
\ml{പോലീസ്} \ml{കള്ളനെ} \ml{ഇന്നലെ} \ml{അറസ്റ്റ്} \ml{ചെയ്തു}.
\end{quote}

\begin{quote}
\emph{pōlīsŭ kaḷḷane innale aṟasṟṟŭ ceytu.}
\end{quote}

\begin{quote}
The police arrested the thief yesterday.
\end{quote}

Past: \ml{അറസ്റ്റ്} \ml{ചെയ്തു}, \emph{aṟasṟṟŭ ceytu}.

The person arrested takes -\ml{എ}: \ml{അവനെ} \ml{അറസ്റ്റ്} \ml{ചെയ്തു}. To be arrested is \ml{അറസ്റ്റിലാകുക}: \ml{അവൻ} \ml{അറസ്റ്റിലായി}, he was arrested.

\begin{cousinweb}
\textbf{\ml{അറസ്റ്റ്} \ml{ചെയ്യുക}} (\emph{aṟasṟṟŭ ceyyuka}) joins \emph{arrest}, borrowed from English, and \ml{ചെയ്യുക}, \textquotedblleft{}to do\textquotedblright{}, the verb that turns a borrowed noun into a Malayalam verb.
\end{cousinweb}

\begin{grammarlens}[title={What you've built}]
One more word for telling what happened.
\end{grammarlens}

\subsection*{Your turn}
\begin{itemize}
\raggedright
  \item \emph{Say it:} \emph{aṟasṟṟŭ ceyyuka}
  \item \emph{Say it:} \emph{aṟasṟṟŭ ceyyuka}, once more
  \item \emph{Say it:} \emph{pōlīsŭ kaḷḷane innale aṟasṟṟŭ ceytu.} — The police arrested the thief yesterday.
  \item \emph{Recall:} say the Malayalam for a criminal, then the Malayalam for an investigation, then say \emph{aṟasṟṟŭ ceyyuka} again
\end{itemize}

\begin{tcolorbox}[breakable,colback=teal!4,colframe=teal!35!black,title={Before you move on}]
What does \ml{അറസ്റ്റ്} \ml{ചെയ്യുക} mean? (\textquotedblleft{}To arrest\textquotedblright{}.) Say it once more.
\end{tcolorbox}
\section[\texorpdfstring{\ml{പിടിച്ചെടുക്കുക} (piṭicceṭukkuka)}{പിടിച്ചെടുക്കുക (piṭicceṭukkuka)}]{\ml{പിടിച്ചെടുക്കുക} (piṭicceṭukkuka) — to seize; to confiscate}

\label{lesson:ML-C365-piticcetukkuka}



\begin{quote}
Before the new one: say the Malayalam for an investigation, then the Malayalam for to arrest.
\end{quote}

\subsection*{You'll want to know: \ml{പിടിച്ചെടുക്കുക}}
\textbf{\ml{പിടിച്ചെടുക്കുക}} — \emph{piṭicceṭukkuka} — \textquotedblleft{}to seize; to confiscate\textquotedblright{}.

\begin{quote}
\ml{പോലീസ്} \ml{കള്ളന്റെ} \ml{ബൈക്ക്} \ml{പിടിച്ചെടുത്തു}.
\end{quote}

\begin{quote}
\emph{pōlīsŭ kaḷḷanṟe baikkŭ piṭicceṭuttu.}
\end{quote}

\begin{quote}
The police seized the thief's motorbike.
\end{quote}

Past: \ml{പിടിച്ചെടുത്തു}, \emph{piṭicceṭuttu}.

\ml{പിടിച്ച്}, having caught, + \ml{എടുക്കുക}, to take: to take away by authority.

\begin{cousinweb}
\textbf{\ml{പിടിച്ചെടുക്കുക}} (\emph{piṭicceṭukkuka}) is a native Dravidian compound of \ml{പിടിക്കുക} (\emph{piṭikkuka}), \textquotedblleft{}to catch, to seize\textquotedblright{} (DEDR 4148), and \ml{എടുക്കുക} (\emph{eṭukkuka}), \textquotedblleft{}to take up\textquotedblright{} (DEDR 851).
\end{cousinweb}

\begin{grammarlens}[title={What you've built}]
One more word for telling what happened.
\end{grammarlens}

\subsection*{Your turn}
\begin{itemize}
\raggedright
  \item \emph{Say it:} \emph{piṭicceṭukkuka}
  \item \emph{Say it:} \emph{piṭicceṭukkuka}, once more
  \item \emph{Say it:} \emph{pōlīsŭ kaḷḷanṟe baikkŭ piṭicceṭuttu.} — The police seized the thief's motorbike.
  \item \emph{Recall:} say the Malayalam for an investigation, then the Malayalam for to arrest, then say \emph{piṭicceṭukkuka} again
\end{itemize}

\begin{tcolorbox}[breakable,colback=teal!4,colframe=teal!35!black,title={Before you move on}]
What does \ml{പിടിച്ചെടുക്കുക} mean? (\textquotedblleft{}To seize; to confiscate\textquotedblright{}.) Say it once more.
\end{tcolorbox}
